\documentclass[11pt]{article}
\usepackage{amsmath,amssymb,amsthm,mathtools,fullpage}
\newtheorem{theorem}{Theorem}
\newtheorem{lemma}[theorem]{Lemma}
\newtheorem{proposition}[theorem]{Proposition}
\newtheorem{corollary}[theorem]{Corollary}
\theoremstyle{definition}
\newtheorem{remark}[theorem]{Remark}
\newtheorem{example}[theorem]{Example}
\newcommand{\Ck}{\mathcal{C}}
\newcommand{\Z}{\mathbb{Z}}
\newcommand{\Q}{\mathbb{Q}}
\newcommand{\vac}{\varnothing}
\DeclareMathOperator{\sgn}{sgn}

\title{A matrix entry equal to $1+t$:\\ the literal gcd of a nested ribbon commutator}
\author{Clio}
\date{5 September 2026}

\begin{document}
\maketitle

\begin{abstract}
Let $R_e(t)$ be the height-graded $e$-ribbon operator on Fock space and, for
$e_1,\dots,e_k\ge2$, let $\Ck_k=[R_{e_1},[R_{e_2},[\dots,[R_{e_{k-1}},R_{e_k}]\dots]]]$.
It is known that $1+t$ divides every matrix entry of $\Ck_k$ (Q81), and that the
$(1+t)$-adic valuation of the gcd of the entries is exactly $1$ if and only if
$e_{k-1}\ne e_k$ (Q83).  The literal statement $\gcd(\Ck_k)=1+t$ was left open,
because it additionally requires that no other irreducible --- in particular
neither $t$ nor $t-1$ --- divides every entry, and the entire hook family used in
Q83 is divisible by $t$ for structural reasons.

We close the gap by exhibiting a \emph{single} matrix entry which is equal to
$1+t$ on the nose.  Put $E=\sum_ie_i$ and $e_{\min}=\min(e_{k-1},e_k)$.  Then
\[
  \bigl\langle (E-1,\,e_{\min})\bigm|\Ck_k\bigm|(e_{\min}-1)\bigr\rangle
  \;=\;\sgn(e_k-e_{k-1})\,(1+t),
\]
for every $k\ge2$ and all $e_1,\dots,e_k\ge2$, with no distinctness hypothesis
beyond $e_{k-1}\ne e_k$.  Consequently $\gcd(\Ck_k)=1+t$ literally, in $\Z[t]$ as
well as in $\Q[t]$ (the entry is primitive, so no content survives), and Q83's
sharpness theorem follows without the closed form for the hook entries.

The proof does not use an involution.  The source state $(e_{\min}-1)$ is a
one-row partition, chosen so that exactly two beads of its Maya diagram can move,
and so that the gap between them equals the displacement of the lower one.  The
paths then split into two sectors --- one in which the beads do not cross, one in
which the lower bead overtakes the upper --- carrying weights $t^0$ and $t^1$
respectively.  The illegal configurations of the two sectors are counted by the
\emph{same} signed sum $X$, so the entry is $-X(1+t)$: the factor $1+t$ is the
statement that a collision forbidden in one sector is exactly a collision
forbidden in the other.  A short sign computation gives $X=-\sgn(e_k-e_{k-1})$.
All statements are verified against an independent bead-model engine.
\end{abstract}

\section{Setting}

Fix an indeterminate $t$.  For a partition $\lambda$ write $|\lambda\rangle$ for the
corresponding basis vector of $\Lambda=\Q(t)\otimes\Lambda_\Z$, and let
\[
  R_e(t)\,|\lambda\rangle=\sum_{\mu}t^{\,h(\mu/\lambda)}|\mu\rangle,
\]
the sum over $\mu\supseteq\lambda$ with $\mu/\lambda$ a connected border strip of size
$e$, and $h$ its height (rows minus one).

We use the Maya (bead) model throughout, in the form fixed in
\texttt{probes/2026-09-04-Q76/abacus.py}: a partition $\lambda$ is encoded by the
bead set
\[
  B(\lambda)=\{\lambda_i-i\ :\ i\ge1\}\subseteq\Z ,
\]
a subset of $\Z$ which contains all sufficiently negative integers and no
sufficiently large one.  Applying $R_e$ means choosing $b\in B$ with $b+e\notin B$
and replacing $B$ by $(B\setminus\{b\})\cup\{b+e\}$; the weight of that hop is
$t^{h}$ with
\[
  h=\#\bigl(B\cap(b,b+e)\bigr),
\]
the number of beads strictly jumped over.  We call such a step a \emph{hop of size
$e$}, and we say the hop \emph{crosses} a bead $c$ if $b<c<b+e$.  A hop never
decreases a bead's position, and beads are indistinguishable only up to the
bookkeeping we now fix: we track individual beads, so that a hop moves a
particular bead.

Throughout, $e_1,\dots,e_k\ge2$, $E:=\sum_ie_i$, and
\[
  \Ck_k=\bigl[R_{e_1},\bigl[R_{e_2},\bigl[\dots,[R_{e_{k-1}},R_{e_k}]\dots\bigr]\bigr]\bigr].
\]
We use two results from the earlier papers.

\begin{theorem}[Q81, \texttt{thm:div}]\label{thm:div}
$1+t$ divides every matrix entry of $\Ck_k$, for every $k\ge2$ and all $e_i\ge2$,
with no distinctness hypothesis.
\end{theorem}

\begin{lemma}[Q81, \texttt{lem:sign}]\label{lem:sign}
Expand $[A_1,[A_2,[\dots,[A_{k-1},A_k]\dots]]]$ into words.  The words that occur are
exactly those whose time-ordered index sequence $w=(w_1,\dots,w_k)$ --- rightmost
operator first, so $w_1$ is the hop performed first --- is unimodal with peak $k$:
$w_1<\dots<w_q=k>w_{q+1}>\dots>w_k$.  Each occurs once, with sign
$\varepsilon(w)=(-1)^{q-1}$.
\end{lemma}

We write $U\subseteq S_k$ for the set of such words.  A word $w\in U$ is determined by
the set $T\subseteq\{1,\dots,k-1\}$ of letters preceding $k$:
\begin{equation}\label{eq:rhoT}
  w=\rho_T=\bigl(\,T\uparrow,\ k,\ D\downarrow\,\bigr),\qquad D:=\{1,\dots,k-1\}\setminus T,
  \qquad \varepsilon(\rho_T)=(-1)^{|T|},
\end{equation}
where $T\uparrow$ means the elements of $T$ in increasing order and $D\downarrow$ those
of $D$ in decreasing order.  In particular
\begin{equation}\label{eq:signsumzero}
  \sum_{w\in U}\varepsilon(w)=\sum_{T\subseteq\{1,\dots,k-1\}}(-1)^{|T|}=0
  \qquad(k\ge2).
\end{equation}

\section{The witness state and its localisation}

\begin{quote}
\emph{Standing notation for the rest of the paper.}  Assume $e_{k-1}\ne e_k$ and put
\[
  e_{\min}:=\min(e_{k-1},e_k),\qquad p:=e_{\min}-2\ \ (\ge0),\qquad
  \lambda^\star:=(e_{\min}-1),\qquad \mu^\star:=(E-1,\,e_{\min}).
\]
\end{quote}
Thus $\lambda^\star$ is a single row and $|\mu^\star|=|\lambda^\star|+E$, as required.
In beads,
\[
  B(\lambda^\star)=\Z_{\le-2}\cup\{p\},\qquad
  B(\mu^\star)=\Z_{\le-3}\cup\{p,\ E-2\} .
\]
(Indeed $\lambda^\star_1-1=e_{\min}-2=p$ and $\lambda^\star_i-i=-i$ for $i\ge2$; and
$\mu^\star_1-1=E-2$, $\mu^\star_2-2=e_{\min}-2=p$, $\mu^\star_i-i=-i$ for $i\ge3$.)
Note $-2<p<E-2$, the last inequality because $E\ge e_{\min}+e_{\max}>e_{\min}$.

Call the bead of $B(\lambda^\star)$ sitting at $p$ the \emph{upper} bead $u$, and the
one at $-2$ the \emph{lower} bead $v$.

\begin{lemma}[localisation]\label{lem:loc}
Let $\pi$ be any sequence of $k$ hops, of sizes $e_{w_1},\dots,e_{w_k}$ in that order,
taking $B(\lambda^\star)$ to $B(\mu^\star)$.  Then every bead of $B(\lambda^\star)$ at a
site $\le-3$ stays put, and no bead ever occupies a site $\le-3$ that it did not
start at.  Hence only $u$ and $v$ move, and at the end $\{u,v\}$ occupies $\{p,E-2\}$.
Exactly one of the following holds.
\begin{enumerate}
\item[\textbf{(NC)}] \emph{Non-crossing sector:} $u$ ends at $E-2$ and $v$ ends at $p$.
  The set $[k]$ of hop labels splits as $S\sqcup\bar S$, the hops of $u$ and of $v$,
  with $\Sigma_S=E-e_{\min}$ and $\Sigma_{\bar S}=e_{\min}$; both parts are nonempty.
\item[\textbf{(C)}] \emph{Crossing sector:} $u$ never moves and $v$ ends at $E-2$,
  performing all $k$ hops.
\end{enumerate}
Here $\Sigma_A:=\sum_{i\in A}e_i$.
\end{lemma}

\begin{proof}
For $y\in\Z$ let $H_y(B):=\#\{x\le y:\ x\notin B\}$ be the number of holes weakly to the
left of $y$; it is finite for every $y$.  A hop moves a bead from some $b$ to $b+e>b$,
so it can only turn a site $\le y$ from occupied to empty (if $b\le y<b+e$), never the
reverse; hence $H_y$ is non-decreasing along $\pi$, and it strictly increases exactly
when a bead crosses from $\le y$ to $>y$.  Now $H_y(B(\lambda^\star))=0=H_y(B(\mu^\star))$
for every $y\le-3$.  Therefore no bead ever moves from a site $\le y$ to a site $>y$,
for any $y\le-3$.  A bead at $x\le-3$ that hopped would move from $x$ to $x+e>x$ and so
cross $y=x\le-3$; so no such bead moves.  Since all hops move beads to the right and
$u,v$ start at $p\ge0$ and $-2$, no bead lands at a site $\le-3$ either.

So only $u$ and $v$ can move, and since $B(\mu^\star)$ differs from $B(\lambda^\star)$
exactly by vacating $-2$ and occupying $E-2$, the pair $\{u,v\}$ finishes on
$\{p,E-2\}$.  The bead $v$ starts at $-2\notin B(\mu^\star)$, so $v$ moves.  If $v$
finishes at $p$ then $u$ finishes at $E-2$, giving (NC): the total displacement of $v$ is
$p-(-2)=e_{\min}$ and that of $u$ is $E-2-p=E-e_{\min}$, and since the $k$ hops carry
total displacement $E$ each hop belongs to exactly one bead, giving the splitting
$S\sqcup\bar S$; $\bar S\ne\varnothing$ as $v$ moves, and $S\ne\varnothing$ because
$\Sigma_S=E-e_{\min}>0$.  Otherwise $v$ finishes at $E-2$ and $u$ at $p$; as $u$ only
moves right and $p$ is its start, $u$ never moves, giving (C).
\end{proof}

\section{The two sectors}

Throughout this section $w\in U$ is a fixed word; $w_s$ is the label of the hop
performed at time $s$.

\subsection{The crossing sector}

\begin{lemma}\label{lem:C}
In sector (C) every legal path has weight exactly $t$, and a word $w\in U$ admits a
(necessarily unique) path in sector (C) if and only if no partial sum
$e_{w_1}+\dots+e_{w_s}$ ($1\le s\le k$) equals $e_{\min}$.  Consequently the total
contribution of sector (C) to $\langle\mu^\star|\Ck_k|\lambda^\star\rangle$ is $-t\,X$,
where
\begin{equation}\label{eq:Xdef}
  X:=\sum_{\substack{w\in U\\ \exists s:\ e_{w_1}+\dots+e_{w_s}=e_{\min}}}\varepsilon(w).
\end{equation}
\end{lemma}

\begin{proof}
In sector (C) the bead $v$ visits $x_0=-2<x_1<\dots<x_k=E-2$ with
$x_s=-2+e_{w_1}+\dots+e_{w_s}$, so the path is determined by $w$.  Throughout, the other
beads sit at $p$ and at $\Z_{\le-3}$; the latter are all $<-2\le x_{s-1}$ and so are
never crossed and never landed on.  The hop from $x_{s-1}$ to $x_s$ is legal iff
$x_s\ne p$, and its weight is $t^{[\,x_{s-1}<p<x_s\,]}$.  Since $x_s=p$ is equivalent to
$e_{w_1}+\dots+e_{w_s}=e_{\min}$, legality of the whole path is the stated condition.
Given legality, $p$ is not among the $x_s$, while $x_0=-2<p<E-2=x_k$; so there is
exactly one $s$ with $x_{s-1}<p<x_s$, and the total weight is $t^1$.  Summing the signs
and using \eqref{eq:signsumzero},
\[
  \sum_{\substack{w\in U\\ \text{no partial sum}=e_{\min}}}\varepsilon(w)\cdot t
  \;=\;t\Bigl(\sum_{w\in U}\varepsilon(w)-X\Bigr)\;=\;-tX. \qedhere
\]
\end{proof}

\subsection{The non-crossing sector}

\begin{lemma}\label{lem:NC}
Fix a splitting $[k]=S\sqcup\bar S$ with $\Sigma_{\bar S}=e_{\min}$ and
$\bar S\ne\varnothing$.  In sector (NC), with $\bar S$ the hops of $v$:
\begin{enumerate}
\item[(a)] every hop has weight $t^0$;
\item[(b)] the path is legal if and only if $\bar S$ is \emph{not} a time-prefix of $w$,
  i.e.\ not equal to $\{w_1,\dots,w_{|\bar S|}\}$.
\end{enumerate}
Consequently sector (NC) contributes $-X'$, where
$X':=\sum_{\bar S}N(\bar S)$, the sum over all $\bar S\subseteq[k]$ with
$\Sigma_{\bar S}=e_{\min}$, and
\begin{equation}\label{eq:Ndef}
  N(\bar S):=\sum_{\substack{w\in U\\ \{w_1,\dots,w_{|\bar S|}\}=\bar S}}\varepsilon(w).
\end{equation}
\end{lemma}

\begin{proof}
Write $P$ (resp.\ $Q$) for the sum of the sizes of the hops of $u$ (resp.\ $v$) performed
so far, so that $u$ sits at $p+P$ and $v$ at $-2+Q$.  Since $Q\le\Sigma_{\bar S}=e_{\min}$
we have $-2+Q\le p\le p+P$ at all times: $v$ stays weakly left of $p$ and $u$ weakly
right of it.

(a) A hop of $u$ goes from $p+P$ to $p+P+e$, and every site in $(p+P,p+P+e]$ is empty:
the beads at $\Z_{\le-3}$ and the bead $v$ are all $\le p\le p+P$.  So $u$'s hops are
always legal and have weight $t^0$.  A hop of $v$ goes from $-2+Q$ to $-2+Q'\le p$.  The
only bead that could lie in the open interval $(-2+Q,-2+Q')$ is $u$, at $p+P\ge p\ge-2+Q'$;
so no bead is crossed and the weight is $t^0$.  The hop is legal iff its target is
unoccupied, i.e.\ iff $p+P\ne-2+Q'$.

(b) By (a), the path is legal iff at each hop of $v$, say at time $s$, we have
$p+P_{<s}\ne-2+Q_{\le s}$, where $P_{<s}$ is the sum of $u$'s hop sizes strictly before
time $s$ and $Q_{\le s}$ that of $v$'s up to and including time $s$.  Since
$p=e_{\min}-2=\Sigma_{\bar S}-2$, this reads $Q_{\le s}\ne\Sigma_{\bar S}+P_{<s}$.  Now
$Q_{\le s}\le\Sigma_{\bar S}$ with equality exactly when $s$ is the last hop of $v$, and
$P_{<s}\ge0$; so the equality $Q_{\le s}=\Sigma_{\bar S}+P_{<s}$ can hold only when
$P_{<s}=0$ and $s$ is $v$'s last hop --- that is, exactly when no hop of $u$ precedes the
last hop of $v$, i.e.\ exactly when the labels $\bar S$ occupy the first $|\bar S|$ time
slots.

Summing over $w\in U$ for fixed $\bar S$ and using \eqref{eq:signsumzero} gives
$\sum_{w\in U}\varepsilon(w)-N(\bar S)=-N(\bar S)$; summing over the admissible $\bar S$
gives $-X'$.
\end{proof}

\begin{lemma}\label{lem:XX}
$X'=X$.
\end{lemma}

\begin{proof}
The partial sums $e_{w_1}+\dots+e_{w_s}$ are strictly increasing in $s$ because every
$e_i\ge2>0$.  Hence for a word $w$ there is at most one $s$ with
$e_{w_1}+\dots+e_{w_s}=e_{\min}$, and if there is one, the set
$\bar S=\{w_1,\dots,w_s\}$ is the unique subset with $\Sigma_{\bar S}=e_{\min}$ which is a
time-prefix of $w$.  So the words counted in \eqref{eq:Xdef} are partitioned according to
that $\bar S$, and $X=\sum_{\bar S}N(\bar S)=X'$.
\end{proof}

\section{The prefix sign sum}

\begin{proposition}\label{prop:N}
Let $k\ge2$ and $\varnothing\ne\bar S\subseteq[k]$, and let $N(\bar S)$ be as in
\eqref{eq:Ndef}.  Then
\[
  N(\bar S)=\begin{cases}
    (-1)^{|\bar S|}   & \text{if }k-1\in\bar S\text{ and }k\notin\bar S,\\[2pt]
    (-1)^{|\bar S|-1} & \text{if }k\in\bar S\text{ and }k-1\notin\bar S,\\[2pt]
    0 & \text{if $\bar S$ contains both of $k-1,k$ or neither.}
  \end{cases}
\]
\end{proposition}

\begin{proof}
Write $r=|\bar S|$ and let $w=\rho_T$ as in \eqref{eq:rhoT}, $m=|T|$, $D=[k-1]\setminus T$.
The first $r$ letters of $\rho_T$ are: the $r$ smallest elements of $T$ if $r\le m$;
$T\cup\{k\}$ if $r=m+1$; and $T\cup\{k\}\cup\{\text{the }r-m-1\text{ largest of }D\}$ if
$r>m+1$.

\smallskip
\emph{Case $k\notin\bar S$.}  Only $r\le m$ can occur.  Then
$\{w_1,\dots,w_r\}=\bar S$ says: $\bar S\subseteq T$ and every element of $T\setminus\bar S$
exceeds $\max\bar S$.  So $T=\bar S\sqcup B$ with $B\subseteq(\max\bar S,\,k-1]$, arbitrary,
and $\varepsilon(\rho_T)=(-1)^{r+|B|}$.  Hence
\[
  N(\bar S)=(-1)^{r}\!\!\sum_{B\subseteq(\max\bar S,\,k-1]}\!\!(-1)^{|B|}
  =(-1)^{r}\cdot\bigl[\,(\max\bar S,k-1]=\varnothing\,\bigr]
  =(-1)^{r}\,[\,\max\bar S=k-1\,],
\]
using $\sum_{B\subseteq Y}(-1)^{|B|}=0$ for $Y\ne\varnothing$.  As $k\notin\bar S$,
$\max\bar S=k-1$ iff $k-1\in\bar S$.

\smallskip
\emph{Case $k\in\bar S$.}  Put $\bar S_0=\bar S\setminus\{k\}$ and $r_0=r-1$.

If $r=m+1$ then $T=\bar S_0$, contributing $(-1)^{r_0}$; this is one term and it always
occurs.

If $r>m+1$ then $T\subseteq\bar S_0$ with $T\ne\bar S_0$, and $\bar S_0\setminus T$ must be
the set of the $|\bar S_0\setminus T|$ largest elements of $D=[k-1]\setminus T$.  Since
$D=([k-1]\setminus\bar S_0)\sqcup(\bar S_0\setminus T)$, this says precisely that every
element of $\bar S_0\setminus T$ exceeds every element of $[k-1]\setminus\bar S_0$.  Let
$m^\ast:=\max\bigl([k-1]\setminus\bar S_0\bigr)$, with $m^\ast:=0$ if that set is empty,
and let $Y:=(m^\ast,\,k-1]$.  By the definition of $m^\ast$ we have $Y\subseteq\bar S_0$,
and the condition is $\bar S_0\setminus T\subseteq Y$, i.e.\ $T\supseteq\bar S_0\setminus Y$.
Writing $T=(\bar S_0\setminus Y)\cup Y'$ with $Y'\subsetneq Y$, we get
$|T|=r_0-|Y|+|Y'|$ and
\[
  \sum_{Y'\subsetneq Y}(-1)^{r_0-|Y|+|Y'|}
  =(-1)^{r_0-|Y|}\Bigl(\sum_{Y'\subseteq Y}(-1)^{|Y'|}-(-1)^{|Y|}\Bigr)
  =\begin{cases}-(-1)^{r_0}&Y\ne\varnothing,\\ 0&Y=\varnothing.\end{cases}
\]
Finally $Y=(m^\ast,k-1]$ is empty iff $m^\ast=k-1$ iff $k-1\notin\bar S_0$.  Adding the two
contributions gives $N(\bar S)=(-1)^{r_0}\bigl(1-[\,k-1\in\bar S\,]\bigr)$, which is the
claim since $r_0=r-1$.
\end{proof}

\begin{lemma}\label{lem:X}
$X=-\sgn(e_k-e_{k-1})$.
\end{lemma}

\begin{proof}
By Lemmas~\ref{lem:NC} and~\ref{lem:XX}, $X=\sum_{\bar S}N(\bar S)$ over all
$\bar S\subseteq[k]$ with $\Sigma_{\bar S}=e_{\min}$.  By Proposition~\ref{prop:N} only
those $\bar S$ containing exactly one of $k-1,k$ contribute.

Suppose $e_{k-1}<e_k$, so $e_{\min}=e_{k-1}$.  If $k\in\bar S$ then
$\Sigma_{\bar S}\ge e_k>e_{k-1}=e_{\min}$, impossible.  If $k-1\in\bar S$ and
$k\notin\bar S$ then $\bar S=\{k-1\}\cup A$ with $A\subseteq[k-2]$ and
$\Sigma_A=e_{\min}-e_{k-1}=0$, forcing $A=\varnothing$ since every $e_i\ge2$.  So the only
contributing set is $\bar S=\{k-1\}$, and $X=(-1)^{1}=-1=-\sgn(e_k-e_{k-1})$.

Suppose $e_{k-1}>e_k$, so $e_{\min}=e_k$.  Symmetrically, $k-1\in\bar S$ would force
$\Sigma_{\bar S}\ge e_{k-1}>e_k$; and $k\in\bar S$, $k-1\notin\bar S$ forces
$\bar S=\{k\}$.  So $X=(-1)^{1-1}=+1=-\sgn(e_k-e_{k-1})$.
\end{proof}

\section{The theorem}

\begin{theorem}[the unit witness]\label{thm:main}
Let $k\ge2$ and $e_1,\dots,e_k\ge2$ with $e_{k-1}\ne e_k$.  Put $E=\sum_ie_i$ and
$e_{\min}=\min(e_{k-1},e_k)$.  Then
\[
  \boxed{\ \bigl\langle (E-1,\,e_{\min})\bigm|\Ck_k\bigm|(e_{\min}-1)\bigr\rangle
  \;=\;\sgn(e_k-e_{k-1})\,\bigl(1+t\bigr).\ }
\]
No distinctness is assumed among $e_1,\dots,e_{k-2}$, and no condition is placed on the
position of the largest size.
\end{theorem}

\begin{proof}
By Lemma~\ref{lem:loc} the entry is the sum of the contributions of sectors (NC) and (C).
By Lemma~\ref{lem:NC} the first is $-X'$, by Lemma~\ref{lem:XX} $X'=X$, and by
Lemma~\ref{lem:C} the second is $-tX$.  So the entry equals $-X(1+t)$, and
Lemma~\ref{lem:X} gives $-X=\sgn(e_k-e_{k-1})$.
\end{proof}

\begin{corollary}[Q85, the literal gcd]\label{cor:gcd}
Under the hypotheses of Theorem~\ref{thm:main}, the greatest common divisor of the matrix
entries of $\Ck_k$ is $1+t$: in $\Q[t]$ up to a unit, and in $\Z[t]$ up to sign.  If
$e_{k-1}=e_k$ then $\Ck_k=0$ and the gcd is undefined (or $0$).
\end{corollary}

\begin{proof}
By Theorem~\ref{thm:div}, $1+t$ divides every entry, so $1+t$ divides the gcd $g$.  By
Theorem~\ref{thm:main} one entry equals $\pm(1+t)$, so $g$ divides $\pm(1+t)$.  Hence
$g=\pm(1+t)$.  The argument is the same in $\Z[t]$ and in $\Q[t]$; in particular no
integer content can occur, because $\pm(1+t)$ is primitive.  For $e_{k-1}=e_k$ we have
$[R_{e_{k-1}},R_{e_k}]=0$ and hence $\Ck_k=0$.
\end{proof}

\begin{corollary}[Q83 sharpness, re-proved]\label{cor:sharp}
For $k\ge2$ and $e_i\ge2$, the $(1+t)$-adic valuation of $\gcd(\Ck_k)$ is exactly $1$ if
and only if $e_{k-1}\ne e_k$.
\end{corollary}

\begin{proof}
Immediate from Corollary~\ref{cor:gcd} and $\Ck_k=0$ when $e_{k-1}=e_k$.
\end{proof}

\begin{remark}[what the mechanism is]
The factor $1+t$ is produced here without an involution and without the toggling
argument of Q83.  It arises because the \emph{same} set of words is forbidden in both
sectors: in (NC) a word fails exactly when $v$'s last hop would land on $u$'s seat $p$
before $u$ has vacated it, and in (C) a word fails exactly when $v$ would land on $p$ at
all.  Lemma~\ref{lem:XX} says these two exclusion sets have the same signed count $X$,
because a landing on $p$ is the same arithmetic event (a partial sum equal to
$e_{\min}$) in both.  Since (NC) carries weight $t^0$ and (C) carries weight $t^1$, the
entry is $-X(1+t)$.  In other words: \emph{the collision that blocks the non-crossing
sector is the collision that blocks the crossing sector}, and $1+t$ records the two ways
the two beads may be arranged around it.
\end{remark}

\begin{remark}[why the source state had to leave the vacuum]
Q83 (\S Gaps, item 1) observes that the hook family
$\langle(E-j,1^j)|\Ck_k|\vac\rangle$ cannot witness $t\nmid\gcd$: every such entry equals
$\sgn(e_k-e_{k-1})(1+t)\sum_T t^{\,j-1-|T|}$ with $j-1-|T|\ge1$, so $t$ divides all of
them.  That is a property of the hook family, not of $\Ck_k$; the correct move is to
change the \emph{source} state rather than to hunt for a better target.  The single row
$\lambda^\star=(e_{\min}-1)$ is the shortest source for which the two-bead gap equals the
lower bead's total displacement, which is exactly what makes Lemma~\ref{lem:NC}(b) a
prefix condition.
\end{remark}

\begin{remark}[the planned second family is a kernel]\label{rem:kernel}
Q83 (\S Gaps, item 3) proposed closing the literal gcd with a second witness family at
$\mu=(2,2,1^{E-4})$, whose entries the Maya engine returns as $\pm t^{a}(t^2-1)$.  Those entries
are as recorded, but $a$ is never $0$, so that family could not have closed the gap either.  The
reason is a symmetry.  Conjugation of partitions reverses the height of a border strip, since
height plus width equals $e-1$; writing $\omega|\lambda\rangle=|\lambda'\rangle$ this says
$\omega R_e(t)\omega=t^{\,e-1}R_e(t^{-1})$, whence
\[
  \bigl\langle\mu'\bigm|\Ck_k(t)\bigm|\vac\bigr\rangle
  =t^{\,E-k}\,\bigl\langle\mu\bigm|\Ck_k(t^{-1})\bigm|\vac\bigr\rangle .
\]
Now $(2,2,1^{E-4})$ is the conjugate of $(E-2,2)$, whose entry is $\pm(t^2-1)$, of degree $2$;
so the entry at $(2,2,1^{E-4})$ is $\pm t^{\,E-k-2}(t^2-1)$, divisible by $t$ whenever
$k\le E-3$.  More generally the displayed symmetry turns ``$t$ divides every entry'' into
``every entry has degree $<E-k$'', so \emph{no family of vacuum entries closed under conjugation
can settle the question by itself}: one must either break the symmetry or, as here, leave the
vacuum.
\end{remark}

\begin{remark}[the $j=e_{\min}$ hook, and a second proof of the corollary]\label{rem:hook}
Specialising the Q83 closed form to $j=e_{\min}$ gives a second, independent route.  The
window $\Sigma_T+e_{\min}\le j<\Sigma_T+e_{\max}$ at $j=e_{\min}$ admits $T=\varnothing$
and nothing else (a nonempty $T$ needs $\Sigma_T\le j-e_{\min}=0$, impossible for
$e_i\ge2$), so
\[
  \bigl\langle (E-e_{\min},1^{e_{\min}})\bigm|\Ck_k\bigm|\vac\bigr\rangle
  =\sgn(e_k-e_{k-1})\,t^{\,e_{\min}-1}(1+t),
\]
a monomial times $1+t$, for \emph{every} admissible tuple.  Any common divisor of the
entries therefore divides $t^{\,e_{\min}-1}(1+t)$, whose divisors in $\Z[t]$ are
$\pm t^{\,i}(1+t)^{\epsilon}$.  So Corollary~\ref{cor:gcd} already follows from
Theorem~\ref{thm:div} together with the weaker statement that \emph{some} entry has
nonzero constant term --- which is the $t=0$ specialisation of Theorem~\ref{thm:main},
i.e.\ sector (NC) alone.  We record this because it isolates what is really needed, and
because the pair of entries
$\bigl(\pm t^{e_{\min}-1}(1+t),\ \pm(1+t)\bigr)$ has gcd $1+t$ visibly.
\end{remark}

\section{Verification}

All computations are in the repository
\texttt{github.com/clio-vega/work-in-progress}, at
\texttt{work-in-progress@00fd9ed:q85/probes/} (locally
\texttt{/home/clio/projects/probes/2026-09-05-Q85/}).  The engines they call live at
\texttt{work-in-progress@00fd9ed:q76/probes/} and \texttt{:q81/probes/}.
Two engines are used, sharing
no primitive: the \emph{Maya engine} (\texttt{probes/2026-09-04-Q76/abacus.py} together
with \texttt{probes/2026-09-04-Q81/nested.py}, which expands the nested bracket into
signed words and hops beads on $\beta$-sets), and the \emph{two-sector model}
(\texttt{twobead.py}), which implements Lemmas~\ref{lem:C}--\ref{lem:NC} directly as a
sum over $(w,\bar S)$ and never constructs a partition.

\begin{enumerate}
\item \textbf{Proposition~\ref{prop:N}} (\texttt{twobead.py}, check 1): the closed form for
$N(\bar S)$ was compared with brute-force enumeration over $U$ for every nonempty
$\bar S\subseteq[k]$ and every $2\le k\le7$ --- $0$ mismatches.
\item \textbf{Lemma~\ref{lem:NC}} (\texttt{twobead.py}, check 2): the identity
$G(\bar S)=-N(\bar S)$, where $G$ is computed by simulating the bead legality condition
$Q_{\le s}<g+P_{<s}$ with $g=\Sigma_{\bar S}$, was verified for all
$\bar S\subsetneq[k]$, all $2\le k\le6$ and all $\vec e\in\{2,\dots,5\}^k$ ---
$288{,}672$ cases, $0$ mismatches.
\item \textbf{Theorem~\ref{thm:main}} (\texttt{verify.py}, checks A and B): the entry
computed by the Maya engine was compared with $\sgn(e_k-e_{k-1})(1+t)$ and, separately,
with the sector model's prediction $-X(1+t)$, for every tuple with $e_{k-1}\ne e_k$ in
$k=2,3,4$ and $e_i\in\{2,\dots,5\}$, plus a sample at $k=5,6$.
\item \textbf{Negative control} (\texttt{verify.py}, check C): replacing $e_{\min}$ by
$e_{\max}$ in the definitions of $\lambda^\star$ and $\mu^\star$ makes the prediction
fail on a positive proportion of tuples, so the comparison is not blind.
\item \textbf{Corollary~\ref{cor:gcd}} (\texttt{verify.py}, check D): the literal gcd of
the full list of entries of $\Ck_k$ from the sources $\vac$ and $\lambda^\star$ was
computed symbolically and equals $1+t$, with integer content $1$.
\item \textbf{Remark~\ref{rem:kernel}} (\texttt{conjscan.py}): the conjugation symmetry was
checked entrywise for $6$ tuples --- $294$ entries, $0$ failures; and no entry at
$\mu=(2,2,1^{E-4})$ has nonzero constant term, over all $k\le4$ and $e_i\in\{2,3,4\}$ with
$e_{k-1}\ne e_k$.
\item \textbf{Remark~\ref{rem:hook}} (\texttt{phase0.py}): the hook entry at
$j=e_{\min}$ was computed from the Maya engine and equals
$\sgn(e_k-e_{k-1})t^{e_{\min}-1}(1+t)$ in every case tested.
\end{enumerate}

\section{What remains}\label{sec:gaps}

\begin{enumerate}
\item \textbf{The vacuum-restricted gcd.}  Corollary~\ref{cor:gcd} is about all matrix
entries of $\Ck_k$, which is the form in which Q83 states its result.  The variant
``$\gcd$ of the entries of $\Ck_k|\vac\rangle$'' is \emph{not} proved here.  It is true
in every case computed (\texttt{t0scan.py}: for all $k\le4$ and $e_i\in\{2,\dots,5\}$
with $e_{k-1}\ne e_k$, some entry of $\Ck_k|\vac\rangle$ has nonzero constant term), but
the witnessing shape varies with $\vec e$: for $\vec e=(2,4,5,3)$ the two-row shape
$(E-e_{\min},e_{\min})=(11,3)$ gives $0$ at $t=0$, and the witnesses are instead
$(9,5)$ and three three-row shapes.  From the vacuum the gap between the two top beads is
$1$, not $\Sigma_{\bar S}$, so Lemma~\ref{lem:NC}(b) becomes a genuine ballot condition
rather than a prefix condition, and Proposition~\ref{prop:N} no longer applies.  This is
the honest remaining gap.
\item \textbf{Sizes $e_i=1$.}  As in Q83, everything is stated for $e_i\ge2$.  The
hypothesis is used twice and essentially: in Lemma~\ref{lem:XX} (strict monotonicity of
partial sums, which needs $e_i\ge1$) and in Lemma~\ref{lem:X} ($\Sigma_A=0\Rightarrow
A=\varnothing$, which needs $e_i\ge1$).  So in fact $e_i\ge1$ suffices for both, but
$p=e_{\min}-2\ge0$ fails at $e_{\min}=1$ and the localisation lemma would need
re-examination.  Untested.
\item \textbf{The order of $N_e$ (Q84).}  Untouched here; Theorem~\ref{thm:main} makes it
independent of Q83 and Q85 alike.
\end{enumerate}

\end{document}
